%print
\newboolean{solution}
\newboolean{isuuid}
\newboolean{link}
\newboolean{isindication}

\newcounter{num}

\definecolor{solutioncolor}{RGB}{0,100,0}
\definecolor{titrecolor}{RGB}{36,76,146}
\definecolor{modernblue}{RGB}{59,130,246}      % Bleu moderne
\definecolor{modernindigo}{RGB}{99,102,241}    % Indigo moderne
\definecolor{moderngray}{RGB}{107,114,128}     % Gris moderne
\definecolor{lightgray}{RGB}{249,250,251}      % Gris très clair
\definecolor{difficultycolor}{RGB}{255,165,0}  % Orange pour la difficulté

%liens vers des ressources extérieures :
\newcommand{\pathcorrige}{https://openyourmath.org/exercise/}
\newcommand{\pathnotebook}{https://github.com/smaxx73/Exercices/blob/main/notebook/}
\newcommand{\pathexercice}{https://raw.githubusercontent.com/smaxx73/Exercices/main/pdf/pdf_sansol/}

\newtcolorbox{marker}[1][]{enhanced,hbox,before skip=2mm,after skip=3mm,fontupper=\small\sffamily,%
	boxrule=0.4pt,left=1mm,right=4mm,top=1pt,bottom=1pt,colback=blue!15,colframe=yellow!20!black,%
	sharp corners,rounded corners=southeast,arc is angular,arc=3mm,%
	underlay={%
		\path[fill=tcbcolback!80!black] ([yshift=3mm]interior.south east)--++(-0.4,-0.1)--++(0.1,-0.2);
		\path[draw=tcbcolframe,shorten <=-0.05mm,shorten >=-0.05mm] ([yshift=3mm]interior.south east)--++(-0.4,-0.1)--++(0.1,-0.2);
	},%
	drop fuzzy shadow,#1
}

\newtcolorbox{solutionbox}{
	enhanced,
	breakable,
	colback=solutioncolor!30!white,
	colframe=solutioncolor,
	boxsep=5pt,
	frame hidden,
	interior style={top color=solutioncolor!15!white, bottom color=solutioncolor!20!white},
	overlay={
		\node[anchor=north west, inner sep=0pt] at (frame.north west) {\tikz{\node[draw, fill=solutioncolor, inner sep=1pt] {\textcolor{white}\checkmark};}};
	}
}

\newtcolorbox{indicationbox}{
	enhanced,
	breakable,
	colback=yellow!30!white,
	colframe=yellow,
	boxsep=5pt,
	frame hidden,
	interior style={top color=yellow!15!white, bottom color=yellow!20!white},
	overlay={
		\node[anchor=north west, xshift=0pt, yshift=0pt] at (frame.north west) {
			\tikz{
				\node[circle, draw=yellow!80!black, fill=yellow!60!black, 
				minimum size=0.3cm, inner sep=0pt] 
				{\textcolor{yellow}{\scriptsize \textbf{i}}};
			}
		};
	}
}

% Commande pour générer les étoiles de difficulté
\newcommand{\difficultystars}[1]{%
	\ifx#1\empty
	\else
	\ifnum#1>0
	\textcolor{difficultycolor}{%
		\ifnum#1=1 $\star$\fi
		\ifnum#1=2 $\star$ $\star$\fi
		\ifnum#1=3 $\star$ $\star$ $\star$\fi
		\ifnum#1=4 $\star$ $\star$ $\star$ $\star$\fi
		\ifnum#1=5 $\star$ $\star$ $\star$ $\star$ $\star$\fi
	}%
	\fi
	\fi
}

% Version modernisée de titrebox
% Version corrigée de titrebox
\newtcolorbox{titrebox}[3][]{
	enhanced,
	breakable,
	boxrule=0pt,
	colframe=white,
	colback=lightgray,
	arc=8pt,  % Coins arrondis modernes
	top=35pt,  % Plus d'espace en haut
	left=20pt,  % Plus d'espace à gauche pour décaler le contenu
	right=8pt,
	bottom=8pt,
	% Pas d'ombre ni de bordure
	overlay unbroken and first={
		% Une seule ligne continue : verticale puis virage vers la droite jusqu'au badge
		\draw[modernblue!60, line width=2pt, rounded corners=3pt] 
		([xshift=12pt]frame.north west) -- 
		([xshift=12pt]frame.south west) -- 
		([xshift=12pt, yshift=-16pt]frame.north west) -- 
		([xshift=22pt, yshift=-16pt]frame.north west);
		
		% Badge principal aligné avec le titre
		\node[anchor=west, inner sep=0pt] at ([xshift=24pt, yshift=-16pt]frame.north west) {
			\tikz{
				\node[
				rounded corners=4pt,
				inner sep=6pt,
				inner ysep=4pt,
				fill=modernblue,
				draw=modernblue!80,
				line width=0.5pt
				] {
					\footnotesize\textbf{\sffamily{\textcolor{white}{Ex #1}}}
				};
			}
		};
		
		% Titre de l'exercice - sans les étoiles maintenant
		\node[anchor=west, inner sep=0pt, text width=\dimexpr\textwidth - 80pt - 60pt\relax] at ([xshift=80pt, yshift=-16pt]frame.north west) {
			\textcolor{moderngray}{\normalsize\bfseries\sffamily{#2}}
		};
		
		% Étoiles de difficulté sous le titre
		\ifx\ExoDifficulte\empty
		\else
		\node[anchor=west, inner sep=0pt] at ([xshift=80pt, yshift=-28pt]frame.north west) {
			\difficultystars{\ExoDifficulte}
		};
		\fi
		
		% UUID avec style moderne (si présent)
		\ifthenelse{\boolean{isuuid}}{
			\node[anchor=north east, inner sep=0pt] at ([xshift=-8pt, yshift=-12pt]frame.north east) {
				\tikz{
					\node[
					rounded corners=3pt,
					inner sep=3pt,
					fill=white,
					draw=moderngray!40,
					line width=0.5pt
					] {
						\ifthenelse{\boolean{link}}{
							\href{\pathcorrige#3}{\textcolor{moderngray}{\tiny\ttfamily{#3}}}
						}{
							\textcolor{moderngray}{\tiny\ttfamily{#3}}
						}
					};
				}
			};
		}{}
	},
	% Gérer les sauts de page - ligne verticale continue
	overlay middle and last={
		\draw[modernblue!60, line width=2pt] 
		([xshift=12pt]frame.north west) -- 
		([xshift=12pt]frame.south west);
	},
}

% Alternative avec style encore plus moderne (optionnel)
\newtcolorbox{moderntitrebox}[3][]{
	enhanced,
	unbreakable,
	boxrule=0pt,
	colframe=white,
	colback=white,
	arc=12pt,
	top=30pt,
	left=6pt,
	right=6pt,
	bottom=6pt,
	drop fuzzy shadow,
	borderline={1pt}{0pt}{lightgray!50},
	overlay unbroken and first={
		% Bande colorée en haut
		\fill[modernindigo, rounded corners=12pt] 
		(frame.north west) rectangle ([yshift=-4pt]frame.north east);
		
		% Badge avec icône
		\node[anchor=north west, inner sep=0pt] at ([xshift=12pt, yshift=-12pt]frame.north west) {
			\tikz{
				\node[
				circle,
				minimum size=28pt,
				fill=white,
				draw=modernindigo!20,
				line width=1pt
				] {
					\textcolor{modernindigo}{\Large\bfseries\sffamily{#1}}
				};
			}
		};
		
		% Titre élégant avec difficulté
		\node[anchor=north west, inner sep=0pt] at ([xshift=50pt, yshift=-20pt]frame.north west) {
			\textcolor{modernindigo}{\LARGE\bfseries\sffamily{#2}}%
			\ifx\ExoDifficulte\empty
			\else
			\hspace{0.5em}\difficultystars{\ExoDifficulte}%
			\fi
		};
		
		% UUID discret
		\ifthenelse{\boolean{isuuid}}{
			\node[anchor=north east, inner sep=6pt] at ([yshift=-12pt]frame.north east) {
				\textcolor{moderngray!70}{\footnotesize\ttfamily{
						\ifthenelse{\boolean{link}}{
							\href{\pathcorrige#3}{#3}
						}{#3}
				}}
			};
		}{}
	},
}

\newcommand{\titre}[1]{%
	\def\TitreExo{#1}
}

\renewcommand{\difficulte}[1]{%
	\def\ExoDifficulte{#1}
}
\newcommand{\contenu}[1]{%
	\def\Contenu{#1}
}

\newcommand{\code}[1]{#1}
\newcommand{\texte}[1]{#1}

\newcommand{\question}[1]{#1}
\newcommand{\reponse}[1]{
	\ifthenelse{\boolean{solution}} 
	{%
		\begin{solutionbox}
			\begin{footnotesize} #1 \end{footnotesize} 
	\end{solutionbox}}{}
} 
\newcommand{\indication}[1]{
	\ifthenelse{\boolean{isindication}} 
	{%
		\begin{indicationbox}
			\begin{footnotesize} #1 \end{footnotesize} 
	\end{indicationbox}}{}
}

\newcommand{\nextexo}{
	\addtocounter{num}{1}
	\vspace{1em}
}

\newcommand{\insertexobis}[5]{%contenu, solution, uuid, lien solution,numerotation,theme
	\noindent
	\input{../src/#1}
	\setboolean{solution}{#2}
	\setboolean{isuuid}{#3}
	\setboolean{link}{#4}
	
	\begin{titrebox}[#5]
		{\bfseries{\TitreExo}}{#1}
		\Contenu
	\end{titrebox}	
}

\newcommand{\insertexo}[6]{%contenu, solution, uuid, lien solution,numerotation,theme
	\Needspace{4\baselineskip}% small guard only
	\noindent
	\input{\path src/#1}
	\setboolean{solution}{#2}
	\setboolean{isuuid}{#3}
	\setboolean{link}{#4}
	\setboolean{isindication}{#6}
	
	\begin{titrebox}[#5]
		{\bfseries{\TitreExo}}{#1}
		\Contenu
	\end{titrebox}	
}

% Version alternative pour un style encore plus moderne (remplacer titrebox par moderntitrebox)
\newcommand{\insertexomodern}[6]{%contenu, solution, uuid, lien solution,numerotation,theme
	\noindent
	% Initialiser la difficulté à vide avant de charger l'exercice
	\def\ExoDifficulte{}
	\input{\path src/#1}
	\setboolean{solution}{#2}
	\setboolean{isuuid}{#3}
	\setboolean{link}{#4}
	\setboolean{isindication}{#6}
	
	\begin{moderntitrebox}[#5]
		{\bfseries{\TitreExo}}{#1}
		\Contenu
	\end{moderntitrebox}	
}

\newcommand{\listexo}[1]{%liste
	\foreach \ex in #1 {
		\nextexo
		\insertexo{\ex}{\solution}{\isuuid}{\link}{\thenum}{\isindication}
	}
}

\newcommand{\insertnotebook}[1]{%contenu, solution, uuid, lien solution,numerotation
	\href{\pathnotebook#1.ipynb}{Lien vers le notebook} 
}

\newcommand{\colonnes}[3]{
	\setboolean{solution}{#1}
	\ifthenelse{\boolean{solution}} {
		\ifnumcomp{#3}{=}{1}
		{}%
		{\begin{multicols}{#3}}%
		}%
		{%
			\ifnumcomp{#2}{=}{1}
			{}%
			{\begin{multicols}{#2}}%
			}
		}
		
		\newcommand{\fincolonnes}[3]{
			\setboolean{solution}{#1}
			\ifthenelse{\boolean{solution}} {
				\ifnumcomp{#3}{=}{1}
				{}%
				{\end{multicols}}%
		}%
		{%
			\ifnumcomp{#2}{=}{1}
			{}%
			{\end{multicols}}%
	}
}

\newcounter{qrcodenum}

\newcommand{\listeqrcode}[2]{%liste d'exercices, initiatialisation du compteur
	\setcounter{qrcodenum}{#2}
	\noindent
	\foreach \exo in #1{
		\begin{minipage}{0.24\textwidth}
			\centering
			\href{\pathcorrige\exo}{\texttt{Ex \theqrcodenum : \exo \vspace{1em}}}
			{\qrcode{\pathcorrige\exo}}
			\par\vspace{2em}
		\end{minipage}
		\stepcounter{qrcodenum}
	}
}